\documentclass[11pt,a4paper]{article}
\usepackage[margin=24mm]{geometry}
\usepackage{amsmath,amssymb,amsthm}
\setlength{\emergencystretch}{2em}
\newtheorem{lemma}{Lemma}
\newtheorem{theorem}[lemma]{Theorem}
\newcommand{\gen}[1]{\langle#1\rangle}
\title{A quotient that excludes the four-point suffix}
\author{Human Compression project: L044 continuation}
\date{}
\begin{document}\maketitle
\begin{lemma}[A fixed point preserves the quotient intersection property]
\label{quotient}
Let $H=\gen{h_0,\ldots,h_k}$ be a string C-group, $k\ge2$,
and let $\pi:H\longrightarrow\overline H$ be a homomorphism
to a permutation group on $X$. Suppose
\[
 N=\ker\pi\le\gen{h_{k-1},h_k},\qquad
 \pi\text{ is injective on }M=\gen{h_0,\ldots,h_{k-1}}.
\]
Suppose some $b\in X$ is fixed by $\pi(M)$ and moved by
$\pi(h_k)$. Then the ordered projected generators form a
string C-group of rank $k+1$.
\end{lemma}
\begin{proof}
Every projected mark is a nonidentity involution: this follows
from injectivity on $M$ for the first $k$ marks and from the
motion of $b$ for the last. String commutation is preserved.
Write $H_I=\gen{h_i:i\in I}$ and $T=\{k-1,k\}$.
We prove the intersection property for arbitrary index sets $I,J$.

If $T\subseteq I$, then $N\le H_I$. For $x\in H_I$ and
$y\in H_J$ with $\pi(x)=\pi(y)$, the element $yx^{-1}$
lies in $N\le H_I$, hence $y\in H_I\cap H_J$.
The original IP gives
$\pi(H_I)\cap\pi(H_J)=\pi(H_{I\cap J})$.
The same proof applies when $T\subseteq J$. This step uses
containment of the entire kernel, not merely a nontrivial
intersection with it.

If $T\not\subseteq I\cup J$, the original IP gives
\[
 N\cap H_{I\cup J}\le H_T\cap H_{I\cup J}
                  =H_{T\cap(I\cup J)}.
\]
The last subgroup is trivial or generated by one of the two
last marks. Each such mark has nonidentity image, so its
intersection with $N$ is trivial. Thus $\pi$ is injective
on $H_{I\cup J}$, and its intersection property follows
directly from that of $H$.

In the remaining case $T\subseteq I\cup J$ but neither set
contains $T$. Interchanging $I,J$ if necessary, one has
$k\in I$, $k-1\notin I$, $k-1\in J$, $k\notin J$.
All marks in $I\setminus\{k\}$ commute with $h_k$, so any
$x\in H_I$ has the form $x=a h_k^\epsilon$, where
$a\in H_{I\setminus\{k\}}\le M$ and $\epsilon\in\{0,1\}$.
If $\pi(x)=\pi(y)$ for $y\in H_J\le M$, the image $\pi(x)$
fixes $b$. Since $\pi(a)$ fixes $b$ and $\pi(h_k)$ moves it,
one must have $\epsilon=0$. Both $x,y$ now lie in $M$, where
$\pi$ is injective. Thus $x=y\in H_I\cap H_J$, and the
required equality follows again.

This proves IP for all subsets. In particular no projected
mark belongs to the subgroup generated by the other projected
marks: applying IP to those complementary index sets would
make its nonidentity cyclic subgroup trivial. The stated
rank is therefore exactly $k+1$.
\end{proof}

\begin{theorem}[The decorated four-point suffix is impossible]
\label{exclude}
Retain the complete critical-head hypotheses and pure-prefix
theorem, in the contrary range $n=2r$. Assume explicitly, as
the strong-induction hypothesis, that for every $12\le m<n$
the natural alternating group $A_m$ has no string C-marking
of rank greater than $\lfloor(m-1)/2\rfloor$.

Suppose $g\ge5$, $r=g+3$, and $\Omega=C\mathbin{\dot\cup}A$,
where $|A|=4$. Let $P_i=G_{\le i}$ and suppose
$P_{g-1}$ fixes $A\cup\{b\}$ pointwise, $b\in C$, while
$P_g$ is transitive on $C$ and contains its pure natural
alternating group. Suppose the last three actual marks are
\[
\begin{split}
 q=\rho_{g+2}&=(a_0\ a_1)(a'_0\ a'_1),\qquad
 A=\{a_0,a'_0,a_1,a'_1\},\\
 s=\rho_{g+1}&=(b\ a_0)(a_1\ a'_1),\\
 u=\rho_g&=(b\ c)(a_0\ a'_0)(a_1\ a'_1)\delta,
 \qquad c\in C\setminus\{b\},
\end{split}
\]
where $\delta$ is supported on $C\setminus\{b,c\}$.
Then no such ambient string C-group exists.
\end{theorem}
\begin{proof}
Put $e=\rho_{g-1}$ and
$v=(a_0\ a'_0)(a_1\ a'_1)$. Every earlier mark
$\rho_i$, $i\le g-2$, fixes $c$: it commutes with $u$ and
fixes $b$. The generator $e$ must move $c$. Otherwise all
earlier marks would fix $A\cup\{b,c\}$ pointwise and the
last three marks would preserve it, contrary to ambient
transitivity. Write $y=c^e\in C\setminus\{b,c\}$.

The complete acquired seven-point alternating-subgroup lemma
applies to the reversed four marks $q,s,u,e$. For exactness,
its labels $1,2,3,4,5,6$ correspond respectively to
$a_1,a_0,a'_1,a'_0,b,c$, so the three literal permutations
in that lemma become exactly $q,s,u$ above. Its last involution
has the form $e=(c\ y)\epsilon$ and fixes $A\cup\{b\}$,
as required. The lemma permits arbitrary disjoint external
factors $\delta,\epsilon$; no natural alternating-tail
assumption from the separate theorem in that packet is used.
Consequently
\[
 \operatorname{Alt}(A\cup\{b,c,y\})\le
                         G_{\{g-1,g,g+1,g+2\}}.       \tag{1}
\]
In particular the pure even permutation $v$ belongs to this
actual suffix.

On the other hand, the exterior image of $P_g$ on $A$ is
exactly $\gen v$: every earlier mark fixes $A$, and $u|_A=v$.
This exterior image is even. Since every ambient permutation
is even, the image of $P_g$ on $C$ is contained in
$\operatorname{Alt}(C)$. Its pure alternating kernel is
already the whole $\operatorname{Alt}(C)$, by hypothesis.
Multiplying $u$ by the inverse of its pure $C$-restriction
therefore gives $v\in P_g$, and in fact
\[
                      P_g=\operatorname{Alt}(C)\times\gen v.
                                                               \tag{2}
\]
Actual ambient IP applied to (1) now gives
\[
                      v\in P_g\cap G_{\ge g-1}=\gen{e,u}.
                                                               \tag{3}
\]

Project (2) to its action on $C$. The entire kernel is
$N=\gen v$, and (3) puts that kernel in the last two marks
of $P_g$. Projection is injective on $P_{g-1}$, since that
group fixes $A$ pointwise, whereas $v$ does not. Finally
$P_{g-1}$ fixes $b$ and $u$ moves $b$ to $c$.
Lemma~\ref{quotient} therefore proves that the actual image
$\operatorname{Alt}(C)$ has a string C-marking of rank $g+1$.
Every projected mark is nonidentity, and independence is
included in that lemma; no rank is inferred merely by counting
the original marks before projection.

Here $n=2r=2g+6$, so $|C|=n-4=2g+2$.
The assumption $g\ge5$ gives $12\le |C|<n$.
The explicit smaller-degree induction hypothesis applies and
requires rank at most
$\lfloor(|C|-1)/2\rfloor=g$. This contradicts rank $g+1$
and excludes the stated configuration.
\end{proof}

\paragraph{Attribution and remaining scope.}
The seven-point alternating subgroup is an acquired proof, used
with its exact literal substitution; its separate natural-tail
theorem is not invoked. The pure-prefix argument and its disclosed
root parity-fibre contribution remain paid. The new quotient lemma
uses a fixed point to verify every projected intersection, closing
the decorated case left by the small-suffix reduction. Root requested
explicit checks of the induction hypothesis, full kernel containment
and projected rank during review preparation; those checks are
included above. The literal suffix is a theorem hypothesis, so
this packet does not assume an unreviewed small-suffix derivation.
Other prefix-orbit configurations, endpoint cuts, opposite heads,
terminal singleton and the complete L044 recovery remain separate.
No new computation, formal theorem, whole-original or question
credit is claimed.
\end{document}
